\begin{proof}[Proof of \cref{obj:lem:jacobi-binomial}]
\begin{enumerate}
\item Set
\[
y:=\frac{z-1}{2}\in\mathbb R.
\]
For \(v\in\mathbb R\) and \(\nu\in\mathbb N_0\), define the rising factorial by
\[
v^{\overline{\nu}}:=\prod_{k=0}^{\nu-1}(v+k),
\qquad v^{\overline{0}}:=1.
\]
The normalized shifted expansion of the Jacobi polynomial gives
\[
P_j^{(a,b)}(z)
=
\frac{(a+1)^{\overline j}}{j!}
\sum_{\nu=0}^{j}
(-1)^\nu\binom j\nu
\frac{(j+a+b+1)^{\overline\nu}}
     {(a+1)^{\overline\nu}}
\left(\frac{1-z}{2}\right)^\nu.
\]
For \(0\le \nu\le j\), splitting the rising product and using the factorial identity for the ordinary binomial coefficient yields
\[
(a+1)^{\overline j}
=(a+1)^{\overline\nu}
 (a+1+\nu)^{\overline{j-\nu}},
\qquad
\binom j\nu\,\nu!\,(j-\nu)!=j!.
\]
Since \(a+1>0\), every factor in \((a+1)^{\overline\nu}\) is positive, including the empty-product case. Thus cancellation is valid, and
\[
\begin{aligned}
\frac{(a+1)^{\overline j}}{j!}\binom j\nu
\frac{(j+a+b+1)^{\overline\nu}}
     {(a+1)^{\overline\nu}}
&=
\frac{(a+1+\nu)^{\overline{j-\nu}}}{(j-\nu)!}
\frac{(j+a+b+1)^{\overline\nu}}{\nu!}\\
&=\binom{j+a}{j-\nu}\binom{j+a+b+\nu}{\nu}.
\end{aligned}
\]
The last equality follows by reversing the factors in each falling product defining the generalized binomial coefficient. Moreover,
\[
(-1)^\nu\left(\frac{1-z}{2}\right)^\nu=y^\nu.
\]
Consequently,
\[
P_j^{(a,b)}(z)
=\sum_{\nu=0}^{j}
\binom{j+a}{j-\nu}\binom{j+a+b+\nu}{\nu}y^\nu.
\]
% lean: CausalSmith.Stat.NoisydoseWeakdesignTransition.jacobi_coefficient_binomial

\item We establish the polynomial identity
\[
\sum_{i=0}^{j}
\binom{\gamma}{j-i}\binom{\gamma_{\mathrm{aux}}}{i}
\xi^i(\xi+1)^{j-i}
=
\sum_{\nu=0}^{j}
\binom{\gamma}{j-\nu}
\binom{\gamma+\gamma_{\mathrm{aux}}-(j-\nu)}{\nu}\xi^\nu,
\qquad
\gamma,\gamma_{\mathrm{aux}}\in\mathbb R,
\]
where
\[
\xi\ \text{is a polynomial indeterminate over }\mathbb R.
\]
For \(0\le i\le\nu\le j\), splitting the falling product gives
\[
\begin{aligned}
\binom{\gamma}{j-i}\binom{j-i}{j-\nu}
&=
\frac{\prod_{k=0}^{j-i-1}(\gamma-k)}
     {(j-\nu)!\,(\nu-i)!}\\
&=
\binom{\gamma}{j-\nu}
\binom{\gamma-(j-\nu)}{\nu-i}.
\end{aligned}
\]
Ordinary binomial symmetry,
\[
\binom{j-i}{j-\nu}=\binom{j-i}{\nu-i},
\]
therefore gives the same product identity with the latter lower index.
% lean: CausalSmith.Stat.NoisydoseWeakdesignTransition.genBinom_choose_product

The generalized Vandermonde convolution for falling-factorial coefficients states
\[
\binom{\gamma+\gamma_{\mathrm{aux}}}{\nu}
=
\sum_{i=0}^{\nu}\binom{\gamma}{i}\binom{\gamma_{\mathrm{aux}}}{\nu-i},
\qquad \nu\in\mathbb N_0.
\]
% lean: CausalSmith.Stat.NoisydoseWeakdesignTransition.genBinom_vandermonde

For \(0\le\nu\le j\), the terms with \(i>\nu\) contribute zero to the coefficient of \(\xi^\nu\) on the left. Expanding \((\xi+1)^{j-i}\), then applying the product identity and Vandermonde convolution, gives that coefficient as
\[
\begin{aligned}
\sum_{i=0}^{\nu}
\binom{\gamma}{j-i}\binom{\gamma_{\mathrm{aux}}}{i}\binom{j-i}{\nu-i}
&=
\binom{\gamma}{j-\nu}
\sum_{i=0}^{\nu}
\binom{\gamma_{\mathrm{aux}}}{i}\binom{\gamma-(j-\nu)}{\nu-i}\\
&=
\binom{\gamma}{j-\nu}
\binom{\gamma+\gamma_{\mathrm{aux}}-(j-\nu)}{\nu}.
\end{aligned}
\]
For \(\nu>j\), every \(0\le i\le j\) satisfies
\[
\nu-i>j-i,
\qquad
\binom{j-i}{\nu-i}=0.
\]
Hence the coefficient is zero on both sides. Equality of all coefficients proves the polynomial identity.
% lean: CausalSmith.Stat.NoisydoseWeakdesignTransition.genBinom_bernstein_poly

\item Evaluate this identity at \(y\), with parameters \(j+a\) and \(j+b\). For \(0\le\nu\le j\),
\[
(j+a)+(j+b)-(j-\nu)=j+a+b+\nu.
\]
Comparison with the expansion in the first step therefore yields
\[
P_j^{(a,b)}(z)
=
\sum_{i=0}^{j}
\binom{j+a}{j-i}\binom{j+b}{i}
y^i(y+1)^{j-i}.
\]
% lean: CausalSmith.Stat.NoisydoseWeakdesignTransition.stdJacobi_binomial_repr

\item Reverse the finite sum by the bijection \(i\mapsto j-i\) of \(\{0,\ldots,j\}\). This gives
\[
P_j^{(a,b)}(z)
=
\sum_{i=0}^{j}
\binom{j+a}{i}\binom{j+b}{j-i}
y^{j-i}(y+1)^i.
\]
Since
\[
y+1=\frac{z+1}{2},
\qquad
2^{j-i}2^i=2^j
\quad(0\le i\le j),
\]
each power product satisfies
\[
y^{j-i}(y+1)^i
=
\frac{(z-1)^{j-i}(z+1)^i}{2^{j-i}2^i}
=
2^{-j}(z-1)^{j-i}(z+1)^i.
\]
Substitution and extraction of the common factor prove
\[
P_j^{(a,b)}(z)
=
2^{-j}\sum_{i=0}^{j}
\binom{j+a}{i}\binom{j+b}{j-i}
(z-1)^{j-i}(z+1)^i.
\]
All product splittings and sum reversals also apply when \(j=0\), using empty products and zeroth powers equal to \(1\).
% lean: CausalSmith.Stat.NoisydoseWeakdesignTransition.stdJacobi_binomial_repr
\end{enumerate}
\end{proof}
